\begin{tikzpicture}[x=1cm,y=1cm,line width=0.8pt]
  % Resistor.
  \draw (0,7.40)--(0.55,7.40) (2.00,7.40)--(2.55,7.40);
  \draw (0.55,7.25) rectangle (2.00,7.55);
  \node[above=2pt] at (1.28,7.55) {$R$};
  \node[right] at (2.95,7.40) {电阻};

  % The inductor is a sequence of upper semicircles.
  \draw (0,6.43)--(0.38,6.43);
  \draw (0.38,6.43)
    arc[start angle=180,end angle=0,radius=0.18]
    arc[start angle=180,end angle=0,radius=0.18]
    arc[start angle=180,end angle=0,radius=0.18]
    arc[start angle=180,end angle=0,radius=0.18]
    arc[start angle=180,end angle=0,radius=0.18]
    --(2.55,6.43);
  \node[above=2pt] at (1.28,6.61) {$L$};
  \node[right] at (2.95,6.43) {电感};

  % Two separate vertical windings and the coupling arch.
  \draw (0.88,5.69)--(0.88,5.49)
    arc[start angle=90,end angle=270,radius=0.14]
    arc[start angle=90,end angle=270,radius=0.14]
    arc[start angle=90,end angle=270,radius=0.14]
    arc[start angle=90,end angle=270,radius=0.14]
    --(0.88,4.15);
  \draw (1.90,5.69)--(1.90,5.49)
    arc[start angle=90,end angle=-90,radius=0.14]
    arc[start angle=90,end angle=-90,radius=0.14]
    arc[start angle=90,end angle=-90,radius=0.14]
    arc[start angle=90,end angle=-90,radius=0.14]
    --(1.90,4.15);
  \draw (1.05,5.60)..controls (1.25,5.76) and (1.53,5.76)..(1.73,5.60);
  \node[left=4pt] at (0.74,4.91) {$L_1$};
  \node[right=4pt] at (2.04,4.91) {$L_2$};
  \node[right] at (2.95,4.91) {互感};

  % Capacitor plates are separated by a genuine gap.
  \draw (0,3.67)--(1.05,3.67) (1.35,3.67)--(2.55,3.67);
  \draw[line width=2pt] (1.05,3.36)--(1.05,3.98)
    (1.35,3.36)--(1.35,3.98);
  \node[above right=1pt] at (1.35,3.85) {$C$};
  \node[right] at (2.95,3.67) {电容};

  % Source conventions in the original: negative left, positive right.
  \draw (0,2.60)--(1.05,2.60) (1.38,2.60)--(2.55,2.60);
  \draw (1.05,2.42)--(1.05,2.78) (1.38,2.24)--(1.38,2.96);
  \node[above=1pt] at (0.86,2.78) {$-$};
  \node[above=1pt] at (1.67,2.78) {$+$};

  \draw (0,1.65)--(0.84,1.65) (1.66,1.65)--(2.55,1.65);
  \fill (0.84,1.65) circle (2.4pt);
  \fill (1.66,1.65) circle (2.4pt);
  \node at (1.25,1.65) {$\sim$};

  \draw (0,0.60)--(0.84,0.60) (1.70,0.60)--(2.55,0.60);
  \draw (1.27,0.60) circle (0.43);
  \node at (1.27,0.60) {$U(t)$};
  \node[above=1pt] at (0.63,0.83) {$-$};
  \node[above=1pt] at (1.91,0.83) {$+$};

  % Hand-drawn brace; this figure needs no additional TikZ library.
  \draw (2.72,3.06)..controls (2.93,3.06) and (2.94,2.84)..(2.94,2.61)
    --(2.94,2.09)..controls (2.94,1.88) and (3.02,1.83)..(3.13,1.78)
    ..controls (3.02,1.73) and (2.94,1.68)..(2.94,1.47)
    --(2.94,0.77)..controls (2.94,0.49) and (2.91,0.19)..(2.72,0.19);
  \node[right,align=left] at (3.24,1.78) {不变的、简谐的和\\任意变动的电压源};
\end{tikzpicture}
